% =====================================================================
%  TRIADE — édition académique v0.10.0 (manuscrit IMRAD, anglais)
%  Corps du document : compilation Tectonic (XeTeX).
%  Les tableaux de données sont GÉNÉRÉS par scripts/extraire_donnees_v100.py
%  (donnee unique : docs/JOURNAL_TRAVAIL.md + docs/DECISIONS.md) et
%  \input ci-dessous — aucune donnée n'est tapée à la main.
% =====================================================================
\documentclass[11pt,a4paper]{article}

\usepackage{fontspec}
\setmainfont{Liberation Serif}
\setsansfont{Liberation Sans}
\setmonofont{Liberation Mono}
\usepackage[french,english]{babel}

\usepackage{amsmath,amssymb}
\usepackage{geometry}
\geometry{a4paper, top=2.6cm, bottom=2.4cm, left=2.3cm, right=2.3cm}
\usepackage{graphicx}
\usepackage{xcolor}
\definecolor{liaison}{HTML}{3681A6}
\usepackage{booktabs}
\usepackage{tabularx}
\newcolumntype{Y}{>{\raggedright\arraybackslash}X}
\usepackage{enumitem}
\setlist{noitemsep, topsep=2pt, leftmargin=1.4em}
\usepackage{microtype}
\usepackage{url}
\usepackage[numbers,sort&compress]{natbib}
\bibliographystyle{unsrtnat}
\setlength{\bibsep}{2pt}
\usepackage{fancyhdr}
\usepackage[colorlinks=true, linkcolor=liaison, citecolor=liaison,
            urlcolor=liaison, unicode=true,
            pdftitle={TRIADE: Enabling LLM Agents to Write Native Code by Hand Under Signed-Intent Governance (Academic Edition v0.11.0)},
            pdfauthor={Alexis Lopez}]{hyperref}
\pagestyle{fancy}
\fancyhf{}
\fancyhead[L]{\small\itshape TRIADE — academic edition v0.11.0}
\fancyhead[R]{\small\thepage}
\renewcommand{\headrulewidth}{0.3pt}

\newcommand{\code}[1]{\texttt{\small #1}}
\newcommand{\RQ}[1]{\textbf{RQ#1}}

\title{\textbf{TRIADE: Enabling LLM Agents to Write Native Code by Hand\\
Under Signed-Intent Governance\\
\large A longitudinal single-case study over eleven milestones, twelve field
runs, and one sealed application}}
\author{Alexis Lopez\thanks{Independent researcher; corresponding author.
Project repository: \code{alexisfrlopez-debug/triade} (GitHub mirror).}}
\date{11 November 2026 — Academic edition v0.11.0 (strengthened evidence)}

\begin{document}
\maketitle
\thispagestyle{fancy}

% ------------------------------------------------------------------
\begin{abstract}
\noindent Large language model (LLM) agents currently produce native code by
generating high-level source and delegating to a compiler toolchain; the
translation from intent to machine code passes through layers the agent
neither writes nor verifies. We report on TRIADE, a development protocol in
which an LLM agent writes native x86-64, RISC-V, and TLS~1.3 code
\emph{directly as assembly notation}, governed by signed layer-I
specifications, an encrypted and fragmented card that binds intentions to
byte ranges, and an epistemology that separates \emph{oracle} evidence
(differential emulation in the laboratory) from \emph{terrain} evidence
(execution on the developer's real machine). We present a longitudinal
single-case study of the protocol's validation programme: eleven milestones
(M0--M10, 26 October -- 11 November 2026), culminating in ``Loupe'', a real
Windows hex-editor application written entirely under the protocol and
delivered to its user. We quantify four research questions on data drawn
from a 9{,}300-line append-only engineering journal: (RQ1) the
escaped-defect rate of changes that had passed all laboratory controls ---
eleven of twelve field iterations surfaced at least one escaped defect,
each named to a class and a mechanism before any fix; (RQ2) the measured
token cost per validated change --- 51 tokens to locate a block, 719 to
edit one, 1{,}088 to add one, against design targets of 64/3{,}000/6{,}000;
(RQ3) the defect classes structurally invisible to emulation, which we
enumerate with their evidence (loader and PE-format semantics, ABI
shadow-space and SSE alignment, Windows error-code conventions,
harness-embedded assumptions); and (RQ4) the safety of hot-patching sealed
code --- zero post-patch regression across the milestone campaigns, with
signed version monotonicity and an armed anti-rollback check. The study is
single-case (N=1), its oracle is derived from the same specifications it
checks, and no baseline compiler-based replication exists yet; we state
these limits explicitly. We claim as contributions the \emph{integration}
of contract-first notation, signed-manifest governance and oracle/terrain
verification into one pipeline for agent-first development --- not the
underlying cryptographic or testing primitives, which are standard.
\end{abstract}

\begin{otherlanguage}{french}
\begin{abstract}
\noindent Les agents de modèles de langage produisent aujourd'hui du code
natif en générant du source de haut niveau délégué à une chaîne de
compilation ; la traduction de l'intention vers le code machine traverse
des couches que l'agent n'écrit ni ne vérifie. Nous rapportons TRIADE, un
protocole de développement dans lequel un agent écrit du code natif
x86-64, RISC-V et TLS~1.3 \emph{directement en notation d'assemblage},
gouverné par des spécifications signées (couche I), une carte chiffrée et
fragmentée qui relie les intentions aux plages d'octets, et une
épistémologie qui sépare la preuve d'\emph{oracle} (émulation
différentielle au laboratoire) de la preuve de \emph{terrain} (exécution
sur la machine réelle du développeur). Nous présentons une étude de cas
unique longitudinale du programme de validation : onze jalons (M0--M10, du
26 octobre au 11 novembre 2026), culminant avec « Loupe », une véritable
application Windows (éditeur hexadécimal) écrite entièrement sous le
protocole et livrée à son utilisateur. Nous quantifions quatre questions
de recherche sur des données tirées d'un journal d'ingénierie append-only
de 9~300 lignes : (RQ1) le taux de défauts échappés des changements ayant
passé tous les contrôles de laboratoire — onze itérations de terrain sur
douze ont révélé au moins un défaut échappé, chacun nommé par classe et
mécanisme avant toute correction ; (RQ2) le coût mesuré en tokens par
changement validé — 51 tokens pour localiser un bloc, 719 pour l'éditer,
1~088 pour l'ajouter, contre des cibles de conception de 64/3~000/6~000 ;
(RQ3) les classes de défauts structurellement invisibles à l'émulation,
énumérées avec leurs preuves (sémantique du chargeur et du format PE,
espace d'ombre et alignement SSE de l'ABI, conventions de codes d'erreur
Windows, hypothèses incrustées dans le harnais) ; et (RQ4) la sûreté du
patch à chaud de code scellé — zéro régression post-patch sur les
campagnes de jalons, avec monotonie de version signée et anti-rollback
armé. L'étude est un cas unique (N=1), son oracle dérive des mêmes
spécifications qu'il vérifie, et aucune réplication avec compilateur de
référence n'existe encore ; ces limites sont énoncées explicitement. Nous
revendiquons comme contribution l'\emph{intégration} — notation à
contrats d'abord, gouvernance à manifestes signés, vérification
oracle/terrain — en un seul pipeline pour le développement piloté par
agent, et non les primitives cryptographiques ou de test sous-jacentes,
qui sont standards.
\end{abstract}
\end{otherlanguage}

% ------------------------------------------------------------------
\section{Introduction}\label{sec:intro}

An LLM-based agent that develops software today almost always works at the
level of high-level source code, and the mapping from that source to
machine code is delegated to a compiler toolchain. This division of labour
has an epistemological consequence: the agent's \emph{intent} is checked
against source text, while the \emph{bytes that execute} are produced by a
translation chain the agent does not write, does not see, and cannot
re-derive. Bugs in that mapping, and drift between what was specified and
what was produced, become a class of defects that no amount of source-level
review can exclude.

TRIADE (\emph{Triple Mirror: Intention, Bytes, Card}) is a protocol
designed for the opposite position: the agent writes native code
\emph{directly}, as assembly notation assembled by a deterministic,
verifiable assembler path --- there is no optimizing or opaque translation
chain; what the agent writes is, by construction, what executes
(assembly notation is verified by an exact round-trip against the emitted
bytes at every build). The protocol was developed and validated by a
single human developer directing an LLM agent, over eleven milestones
between 26 October and 11 November 2026, and is documented in a
9{,}348-line append-only work journal, a 4{,}421-line append-only decision
log (D1--D181), and sealed artifacts whose SHA-256 fingerprints are
consigned at each sealing ceremony.

This paper is the first academic edition of that effort. Its purpose is
not to argue that hand-written machine code is superior to compiled code
in general; it is to report, with the data the protocol itself produced,
what happens when an LLM agent develops real native software this way ---
what it costs, what breaks, what laboratory verification could not see,
and what governance was needed to keep the process falsifiable. The
protocol's cryptographic and testing primitives (signature schemes,
secret sharing, differential testing, key ceremonies) are standard and
cited; the contribution claimed is their \emph{integration} into one
agent-first pipeline, together with the oracle/terrain epistemology and
the signed-intent governance described below.

The study is structured around four research questions:

\begin{itemize}
\item \RQ{1} \emph{Escaped-defect rate.} When a change passes every
laboratory control (differential oracles, harnesses at three fixed seeds,
structural checks, seal guards), how often does it still fail on the real
target machine, and at what granularity can the failure be named?
\item \RQ{2} \emph{Token cost per validated change.} What does it cost,
in LLM context tokens, to locate, edit, or add a block of sealed native
code under the protocol's card-based indexing, measured rather than
estimated?
\item \RQ{3} \emph{Emulation-invisible defect classes.} Which classes of
defects are structurally invisible to a Unicorn-based emulation harness,
and therefore escape to the terrain regardless of laboratory coverage?
\item \RQ{4} \emph{Hot-patch safety.} Can code be patched in place while
sealed artifacts remain byte-stable, with zero post-patch regression and
an armed rollback barrier, across ISAs?
\end{itemize}

The remainder of the paper describes the protocol and the study methods
(Section~\ref{sec:methods}), the results per research question
(Section~\ref{sec:results}), a discussion (Section~\ref{sec:discussion}),
threats to validity (Section~\ref{sec:threats}), related work
(Section~\ref{sec:related}), and declarations.

% ------------------------------------------------------------------
\section{Methods}\label{sec:methods}

\subsection{The TRIADE protocol in brief}\label{sec:protocole}

TRIADE organises development in three indissociable layers. The
\emph{intention layer} (layer I) holds one normalized specification per
code block: identifier, intent, pre/post-conditions and invariants, named
failure cases, verification vectors, fixtures, provenance and the
decisions that shaped it. The \emph{block layer} (layer B) holds the
assembly notation of each block, at fixed addresses, with deterministic
assembly and a mandatory round-trip check (assemble, disassemble,
re-assemble; any divergence refuses the build). The \emph{card layer}
(layer C) is an encrypted, fragmented, signed index (the \emph{card}) that
binds every block to its byte range, contract and fingerprints; only the
card is needed to locate or edit a block, never the whole binary.

Eleven rules (R1--R11) fix the discipline; the load-bearing ones for this
study are: R2 (\emph{written before code}: every block is specified before
any notation byte exists); R4 (\emph{stable contract}: a change may not
silently alter a block's contract); R5 (\emph{verify then write}: the
agent must re-derive the check before writing); R6 (crystallisation: an
artifact may only be sealed after ten consecutive identical rebuilds);
R7 (\emph{sealed inheritance}: inherited bytes live at their sealed
addresses, byte-identical); R10 (load-time verification: checksums per
block, signed version monotonicity, anti-rollback floor); R11 (no
optimization without measurement).

The development loop itself (update an intention, re-derive the oracle,
edit the notation, rebuild, verify, seal) and the seal ceremony model are
described in the protocol specification v0.9.0, which this edition
accompanies; we summarise only what the evaluation needs.

\paragraph{Governance: signed intent and append-only records.}
Every layer-I tree is condensed and signed with an Ed25519 key kept
strictly offline; a build refuses an unsigned layer I. Structural
decisions are consigned in an append-only decision log --- each entry has
an identifier, date, status, context, the retained choice, rejected
alternatives and amendment links; entries are never rewritten or
back-dated. The eleven signed trees at the time of the study are verified
by a guard that must pass \emph{before} any build (14/14 at the last
run). Sealing a release is a \emph{fact of the card}: the sealed tag
lives in the card's build fragment; the artifact's bytes never move
(Fig.~\ref{tab:empreintes} lists the sealed artifacts and their
fingerprints).

\paragraph{The oracle/terrain epistemology.}
The protocol distinguishes two evidence regimes. \emph{Oracle evidence} is
produced in the laboratory: differential oracles derived from the layer-I
specifications (never exported, always regenerated from I), scenario
oracles that replay expected screens byte-for-byte, structural oracles
that read the composed image instruction by instruction, and Unicorn-based
emulation harnesses with Windows ABI shims. \emph{Terrain evidence} is
produced on the developer's real machine: the sealed artifact, a
scenario-loader script, and witnesses (byte-exact logs, showcase runs)
that are pushed back to an evidence vault and cross-checked (``recoup'')
against the laboratory's expectations. The study's unit of failure is the
\emph{named defect}: a failure is not counted until a class, a mechanism
and an evidence piece have been assigned to it, and every fix is a
documented surgical revision (kit-only or piece-only, with byte diffs).

\paragraph{The agent.}
The developer directed an LLM agent (Z.ai GLM family, ``the laboratory'')
through short natural-language instructions; the agent executed all
laboratory work (notation, oracles, harnesses, builds, ceremonies,
consignation) in a sandboxed Linux environment with no access to the
offline signing key. The complete transcript of the collaboration is the
append-only journal and work log themselves --- every session ends with a
consigned entry listing what was done, what was verified and what remains
--- plus the decision log; these, with the evidence vault, are the
primary sources of this paper. The agent is described here and in the
acknowledgements; it is not an author.

\subsection{The validation programme and the case study}\label{sec:cas}

The protocol was validated by eleven milestones, each with measured exit
criteria (Table~\ref{tab:jalons}); M0 established the artifact chain,
M2 hot patching on a real machine, M3 a million-case fuzzing campaign,
M4 an encrypted Schannel tunnel, M6/M7 portability to Linux and RISC-V,
M8 direct syscalls, M9 pure TLS~1.3 (28 blocks, official vectors from
FIPS~180-4, RFC~2104/4231, RFC~5869, RFC~7748, SP~800-38D, handshake
replayed on the RFC~8448 transcript, and a real encrypted session with a
real server), and M10 the full-scale test: \emph{Loupe}, a real Windows
hex-editor application --- 18 blocks in five families (file, digests,
rendering, navigation, editing) --- written entirely under the protocol,
delivered as a field kit to the developer's machine, and run twelve times
(Table~\ref{tab:runs}). Verification volumes per milestone are
consolidated in Table~\ref{tab:oracles}.

\subsection{Measurements}\label{sec:mesures}

\emph{Token budgets} were instrumented at M1 (decision D31/D34) with the
\code{cl100k\_base} tokenizer over the sealed card and real sources:
tokens to locate a block (read one \code{frag-loc} entry), to edit one
(additional \code{frag-con} + notation), to add one (additional
profile thresholds). \emph{Escaped defects} are counted per terrain
iteration: a delivered kit that passed all laboratory controls and then
failed on the terrain counts one escaped defect per named class.
\emph{Hot-patch safety} is measured as post-patch regressions (target:
zero) and as the byte-stability of sealed artifacts across patch
campaigns. All numbers reported in Section~\ref{sec:results} are
generated programmatically from the journal and decision log by the
extraction script shipped with this edition; the manuscript contains no
hand-typed data value.

% --- Tableaux générés depuis le journal (source unique de vérité) ---
% Fichier GÉNÉRÉ par extraire_donnees_v100.py — ne pas éditer à la main.
% Chaque valeur est extraite/validée depuis docs/JOURNAL_TRAVAIL.md et docs/DECISIONS.md (édition académique v0.10.0, D182).

\begin{table}[htbp]\centering\footnotesize
\caption{Milestones M0--M10 of the TRIADE validation programme (26 Oct -- 11 Nov 2026). Every milestone was closed on measured exit criteria; sources: work journal entries and decision log.}
\label{tab:jalons}
\begin{tabularx}{\textwidth}{@{}ll lX@{}}
\toprule
\textbf{Milestone} & \textbf{Sealed tag} & \textbf{Date} & \textbf{Lab evidence (exit criteria met)} \\
\midrule
\textbf{M0} & \code{0.1.0} & 10-26 & round-trip 100\\%; executable < 4 KiB \\
\textbf{M1} & \code{0.2.0} & 10-26 & oracle 21/21; first crystallisation \\
\textbf{M2} & \code{0.3.0} & 10-27 & conformity 68/68; 8 instrumented iterations \\
\textbf{M3} & \code{0.4.0} & 10-28 & $10^6$ fuzzed cases $\times$2 seeds; 161 conformity controls \\
\textbf{M4} & \code{0.5.0} & 10-29 & 560-case SSPI scenario oracle $\times$4 seeds; 12 mutations caught \\
\textbf{M5} & \code{0.6.0} & 11-04 & per-build fingerprints active \\
\textbf{M6} & \code{0.7.0} & 11-05 & 7/7 conformity suite; layer-I diff zero \\
\textbf{M7} & \code{0.8.0} & 11-05 & round-trip 100 \% on new ISA; fence.i patch \\
\textbf{M8} & \code{0.9.0} & 11-05 & 61 harness greens; 131 socket-socle controls \\
\textbf{M9} & \code{1.0.0} & 11-07 & 28 blocks; 589 structural locks; handshake 54/54 $\times$3 seeds; cross-lock 86/86 \\
\textbf{M10} & \code{loupe 0.1.0 / 0.1.1} & 11-11 & 18 blocks; 159 controls $\times$3 seeds; cross-lock 26/26; seal guard 14/14 \\
\bottomrule
\end{tabularx}\end{table}

\begin{table}[htbp]\centering\footnotesize
\caption{Measured token budgets (M1 instrumentation, tokenizer cl100k\_base) against design targets, and generation-1 lineage volumes (D34).}
\label{tab:budgets}
\begin{tabular}{@{}lrrl@{}}
\toprule
\textbf{Operation} & \textbf{Target} & \textbf{Measured} & \textbf{Source} \\
\midrule
Locate a block & 64 & 51 & D31/D34, M1 measurement (cl100k\_base) \\
Edit a block & 3000 & 719 & D31/D34, M1 measurement \\
Add a block & 6000 & 1,088 & D34, M1 (draft 1086, +2 drift re-measured) \\
Whole lineage (specs) & --- & 8,774 & D34, generation-1 card \\
Whole lineage (drafts) & --- & 3,973 & D34 \\
Whole lineage (delivered notations) & --- & 5,099 & D34 \\
Whole lineage (TOTAL) & --- & 13,873 & D34 \\
\bottomrule
\end{tabular}\end{table}

\begin{table}[htbp]\centering\footnotesize
\caption{Verification instruments and volumes per milestone. All numbers are drawn from the corresponding journal entries.}
\label{tab:oracles}
\begin{tabularx}{\textwidth}{@{}l X@{}}
\toprule
\textbf{Milestone} & \textbf{Instrument (measured volume)} \\
\midrule
\textbf{M1} & mini-libc oracle — 21/21 \emph{(D32; tag 0.2.0)} \\
\textbf{M2} & echo conformity + differential — 68/68; 8 iterations, 8 defect classes \emph{(D43; tag 0.3.0)} \\
\textbf{M3} & fuzz campaign + profile conformity suite — 1,000,000 cases $\times$2 seeds; 161 controls \emph{(D72/D73)} \\
\textbf{M4} & SSPI scenario oracle + per-block differential — 560 cases $\times$4 seeds; 10,000 fuzzed cases $\times$2 seeds per block; 165k leaf differentials \emph{(D68-D93)} \\
\textbf{M5} & build fingerprint registry + verifier — per-build sha256 registers \emph{(D134, D148)} \\
\textbf{M6} & profile conformity suite (Linux) — 7/7; lab recoup 22/22 \emph{(D154)} \\
\textbf{M7} & RISC-V round-trip + cross-profile lock — 100 \% round-trip; 29/29 recoup \emph{(D155)} \\
\textbf{M8} & sys harness on synthetic ntdll fixtures — 61 greens; 131 socle controls; 30/30 recoup \emph{(D160-D165)} \\
\textbf{M9} & tls13 structural oracle + handshake harness + cross-lock — 589 structural locks; handshake 54/54 $\times$3 seeds; cross-lock 86/86; 560+365+131 replayed \emph{(D166-D170)} \\
\textbf{M10} & Loupe: 3 Unicorn harnesses + structural oracle + seal guard — 159 controls $\times$3 seeds (61+53+45); 121 structural locks; cross-lock 26/26; seal guard 14/14; crystallisation 26 x 10/10 \emph{(D172-D181)} \\
\bottomrule
\end{tabularx}\end{table}

\begin{table}[htbp]\centering\scriptsize
\caption{The twelve terrain runs of the Loupe field test (M10). Every failure was \emph{named} (a class, a mechanism, an evidence piece) before being fixed; every fix was documented as a surgical revision. The sealed artifacts never changed after their sealing run. Sources: journal entries 204--222, decisions D178/D181.}
\label{tab:runs}
\begin{tabularx}{\textwidth}{@{}r l X X@{}}
\toprule
\textbf{\#} & \textbf{Kit} & \textbf{Outcome (named defect)} & \textbf{Documented fix} \\
\midrule
1 & \code{dev.1} & loader refusal 0xc000007b. Missing BASERELOC: ASLR image without relocation table (lab harness parses, never loads) & dev.2: minimal 12-byte BASERELOC, 7 bytes in 6 ranges \\
2 & \code{dev.2} & same refusal 0xc000007b. Necessary-not-sufficient: page hole 0x3000-0x4000 + massive BSS outside the proven envelope of 8 prior artifacts & dev.3: image within envelope (VS==RS, no BSS) \\
3 & \code{dev.3} & loader mapped the image (decisive step); T0 witness failed. [CADRE-NE]: .data initializer wrong (sentinel at +8 - a SIZE read as an offset instead of +528) & dev.4: initializer = exact mirror of harness CADRE\_VIERGE; guard healed (control derives from plan, never from the piece it guards) \\
4 & \code{dev.4} & T0 all green; T1 named -6; showcase died instantly. Stack-argument class: msx64 args 5+ at home space instead of [rsp+0x20]+ (loupe-01/05/06) + SSE alignment (loupe-10/11) & dev.5: slots fixed; msx64 ABI audit born at build (static, catches SSE class); falsification control V7 added \\
5 & \code{dev.5} & first run acquitting the piece; T1 named -7 VERROUILLE. [SESSION]: kit itself - scenario-as-loader holds loupe.exe as a live section, Windows refuses gle 32 to open it for write & dev.6 KIT-ONLY (artifact unchanged): byte-copy shadow temoin-session.bin \\
6 & \code{dev.6} & two defects from one witness. False save error + missing showcase thread; [HINSTANCE-PEB]: hInstance read from PEB instead of GetModuleHandleW(NULL) & piece dev.6 (r12 = GetModuleHandleW(NULL)) + kit dev.7 \\
7 & \code{dev.7} & T0-T5 all green (save proven on terrain); then two defects. [EAX-SANS-SIGNE] at kit (unsigned 32-bit reads of negative codes) + [D54-QUATRIEME] at piece (4 contiguous FirstThunk ranges without NULL separators) & piece dev.7 (IAT relocated 0x3100, 4 ranges separated by NULL, 63 call sites repointed) + kit dev.8 \\
8 & \code{dev.8} & T0-T7 all green for the FIRST time (functional contract fully proven). Probe died on its own Python; showcase still mute (5th death) & kit dev.9 (probe made load-bearing) \\
9 & \code{dev.9} & T0-T7 all green second time; 6th showcase death named by system then by lab. [OFT-QWORD]: 4 OFT tables at QWORD stride instead of DWORD - the loader walks 10.18.6.3 entries at DWORD stride; first cause of ALL deaths found on the piece bytes & piece dev.8 (89 bytes in 52 ranges, sealed 0.1.0 candidate) + kit dev.10 \\
10 & \code{dev.11} & T0-T7 all green third time; SHOWCASE ALIVE FIRST TIME (0 and 0, no WER). [FREE-SANS-TYPE] at kit: FreeLibrary called without argtypes, 64-bit HMODULE truncated to C int - ArgumentError killed cascade 205 & kit dev.12 (argtypes=[c\_void\_p], except Exception widened, probe armored try/finally) \\
11 & \code{dev.12} & T0-T7 all green 4th time; cascade 205 complete 1st time; showcase alive 2nd time; green holds entire. [SONDE-EXE] named by lab on pieces: structural false negative - library loader NEVER resolves an EXE's imports; liaison proven by the living showcase & seal loupe 0.1.0 on recouped run 39/39 (D178): notation debt paid (587$\rightarrow$0), crystallisation 26x10/10, anti-rollback armed \\
12 & \code{dev.1} & T0-T7 all green 5th time; cascade 205 2nd time; showcase alive 3rd time (0 and 0, first of piece 0.1.1). [T0-PERIME] residual named (dev.12-era banner text - no gate touched) & seal loupe 0.1.1 (D181): [R10-PERIMEE] debt paid at birth (17 sums refreshed), card generation 17, artifact intact \\
\bottomrule
\end{tabularx}\end{table}

\begin{table}[htbp]\centering\footnotesize
\caption{Sealed artifacts (SHA-256, verified byte for byte against the card and journal at edition time).}
\label{tab:empreintes}
\begin{tabular}{@{}l r l@{}}
\toprule
\textbf{Artifact} & \textbf{Bytes} & \textbf{SHA-256 (first 16 hex)} \\
\midrule
\code{loupe.exe 0.1.1} & 127,488 & \code{c111510535767a40\ldots} \\
\code{loupe.exe 0.1.0} & 127,488 & \code{054de960c31425d0\ldots} \\
\code{proxy.exe 1.0.0} & 65,024 & \code{446279f3ab356900\ldots} \\
\bottomrule
\end{tabular}\end{table}

\begin{table}[htbp]\centering\footnotesize
\caption{Strengthened evidence (v0.11.0, D185): six new instruments, all run with fixed seeds and reproducible. Sources: scripts/preuves\_v110/.}
\label{tab:preuves110}
\begin{tabularx}{\textwidth}{@{}l l X l@{}}
\toprule
\textbf{Instrument} & \textbf{Target} & \textbf{Measured} & \textbf{Detail} \\
\midrule
Mutation testing (16 mutants/bloc, graine  & loupe & 230/288 killed (79.9\%) & 18 blocks \\
Mutation testing (16 mutants/bloc, graine  & tls13 & 397/448 killed (88.6\%) & 28 blocks \\
Mutation testing (16 mutants/bloc, graine  & sys & 67/96 killed (69.8\%) & 6 blocks \\
Constant-time (dudect d'émulation, N=60) & tls13-01..04 (sha-256) & t = 0.0 & CONSTANT \\
Constant-time (dudect d'émulation, N=60) & tls13-05 (hmac-sha256, la clé) & t = 0.0 & CONSTANT \\
Constant-time (dudect d'émulation, N=60) & tls13-06 (hkdf-extract, l'IKM) & t = 0.0 & CONSTANT \\
Constant-time (dudect d'émulation, N=60) & tls13-13/14 ($\times$25519, la clé privée) & t = 0.0 & CONSTANT \\
Constant-time (dudect d'émulation, N=60) & tls13-15..19 (aes-gcm, la clé) & t = 0.0 & CONSTANT \\
Constant-time (dudect d'émulation, N=60) & loupe-06 (md5, le contenu) & t = 0.0 & CONSTANT \\
Coverage-guided fuzzing (graines 20261111/ & loupe-01 & 60.5\% coverage, 0 crashes, 0 canary violations & corpus 18 \\
Coverage-guided fuzzing (graines 20261111/ & http-10 & 358 instructions covered, 0 crashes & corpus 17 \\
Oracle independence (distinct seed and mod & I1 tables MD5 T[64] : plan ↔ dérivation  & 0 divergence(s) & zéro \\
Oracle independence (distinct seed and mod & I1 décalages MD5 S[64] : plan ↔ RFC 1321 & 0 divergence(s) & zéro \\
Oracle independence (distinct seed and mod & I2 vecteurs RFC 1321 : spec (par citatio & 0 divergence(s) & 7 vecteurs : harnais == hashlib en clair \\
Oracle independence (distinct seed and mod & I3 différentiel sha256 : 200 messages os & 0 divergence(s) & zéro écart \\
Oracle independence (distinct seed and mod & I3 différentiel md5 : 200 messages os.ur & 0 divergence(s) & zéro écart \\
Oracle independence (distinct seed and mod & I4 K[64] SHA-256 : plan ↔ racines cubiqu & 0 divergence(s) & zéro \\
Bounded static memory safety & 52 blocks & 1181 stack accesses verified, 0 out-of-frame & 463 unresolved (named) \\
Comparative benchmarks (same emulator) & sha256(1024 o) & 69132 ins/op & 5436 bytes \\
Comparative benchmarks (same emulator) & sha256(1024 o) & 77505 ins/op & 3072 bytes \\
Comparative benchmarks (same emulator) & sha256(1024 o) & 4524 ins/op & 765 bytes \\
\bottomrule
\end{tabularx}\end{table}


% ------------------------------------------------------------------
\section{Results}\label{sec:results}

\subsection{RQ1: escaped-defect rate}\label{sec:rq1}

Table~\ref{tab:runs} reports the twelve terrain runs of the Loupe case
study. Eleven of the twelve iterations surfaced at least one escaped
defect --- a defect that had passed every laboratory control in force at
the time, including the three Unicorn harnesses (159 controls at three
fixed seeds), the structural oracle (43 then 121 locks) and the seal
guard. The twelfth run (the 0.1.1 artifact, whose card-debt had been paid
at birth) was green end-to-end. Counting named classes across the
campaign: eleven named defect classes over eleven failing iterations,
each with a named mechanism and an archived evidence piece
(\code{preuves-run-loupe-*} in the evidence vault). The showcase ---
the application simply being \emph{used} on the terrain --- was alive with
zero error-reporter events at runs 10, 11 and 12, after six consecutive
deaths.

The earlier milestones show the same shape at smaller scale: the M2
instrumentation campaign took eight iterations and named eight classes
before the hot-patch loop was proven on a real machine (D43); the M3
proxy needed eight terrain deliveries before the first complete run
(D63), and its killer defects (kernel32 shadow-space violations, a
\code{sub rsp,0x7100} without stack probing) were all of classes the
emulator could not see. We treat this as the study's central empirical
finding: \emph{passing a complete laboratory programme does not transfer
to the terrain}, and the protocol's response is not more laboratory
coverage but a terrain evidence regime in which every escape becomes a
named, falsified class --- each of the eleven Loupe classes is now
guarded by a control or an audit born from its own mechanism (e.g.\ the
msx64 ABI audit born at build from the SSE class of run 4).

\subsection{RQ2: token cost per validated change}\label{sec:rq2}

Table~\ref{tab:budgets} shows the instrumented budgets. Locating a block
costs 51 tokens measured against a 64-token design target; editing a
block 719 against 3{,}000; adding a block 1{,}088 against 6{,}000. The
whole generation-1 lineage --- all specifications, drafts and delivered
notations --- totalled 13{,}873 tokens (D34). Under the protocol's
card-fragment indexing, no operation of the update loop ever requires
reading the binary; the card is the only index, and full-card audits are
exceptional, logged operations. These figures predate the Loupe
application and were not re-measured for it; we report them as the
protocol's M1 instrumentation result (a limitation is discussed in
Section~\ref{sec:threats}).

\subsection{RQ3: defect classes invisible to emulation}\label{sec:rq3}

The eleven named classes of the Loupe campaign group into five families;
each was invisible to the emulation harness for a structural reason, not
a coverage gap:

\begin{enumerate}
\item \emph{Loader and image semantics} (runs 1--3): an ASLR image
without a relocation table (refusal \code{0xc000007b}); an image outside
the proven envelope of eight prior artifacts (page hole, massive BSS);
an uninitialized \code{.data} session frame (the initializer existed only
in the harness's ideal memory). The harness parses images; it never loads
them.
\item \emph{ABI conventions with hardware consequences} (run 4): msx64
stack arguments 5+ written at the home space instead of
\code{[rsp+0x20]+}, and stack misalignment at call sites
($\mathrm{rsp} \equiv 8 \bmod 16$) --- harmless under emulation, fatal
with a real SSE-using \code{gdi32}. This class gave birth to a static
msx64 ABI audit at build time, which then caught two further alignment
defects no one had seen.
\item \emph{Operating-system loader behaviour} (runs 5, 9, 11): Windows
refusing a write-open on the file of a live mapped section
(\code{gle 32}); import tables walked by the loader at \code{DWORD}
stride while the build wrote \code{QWORD} entries --- the first cause of
all six showcase deaths, found on the artifact's bytes; and the
structural false negative that a library loader never resolves an EXE's
imports (the ``liaison'' probe could never succeed as designed --- the
showcase itself is the proof).
\item \emph{Error-code conventions} (runs 6--7): an \code{hInstance}
read from the PEB instead of \code{GetModuleHandleW(NULL)}; negative
return codes read as unsigned 32-bit values by the kit while the piece
delivered them sign-extended.
\item \emph{Harness- and kit-embedded assumptions} (runs 8, 10 and the
pre-run fix of run 11): a probe killed by its own unguarded Python; a
64-bit \code{HMODULE} truncated by \code{ctypes} to a C \code{int}
because \code{argtypes} was not set; a bare \code{ret} neutralizing an
entry point whose return value the loader reads in \code{DllMain}
semantics. These are defects of the \emph{test rig}, named with the same
discipline as defects of the artifact.
\end{enumerate}

The pattern across families: the emulator verifies \emph{logic} well and
\emph{environment} not at all. The protocol's answer is asymmetric ---
add the static audits that can see the environment-sensitive classes
(ABI, image format), and route the rest through terrain evidence with
named-defect accounting, rather than pretending the harness can be
completed.

\subsection{RQ4: hot-patch safety}\label{sec:rq4}

The M2 campaign demonstrated the full loop on a real machine: replacing
the echo block \emph{while executing}, verifying by comparison,
rolling back, re-verifying --- with zero post-patch regression across the
eight instrumented iterations (D43). On RISC-V (M7), where fixed-width
instructions forbid the x86-64 relative trampoline, the patch path uses
a scratch register plus \code{fence.i}, validated by the same oracle
programme (29/29 recoup on crossed pieces). Load-time verification (R10)
checks per-block checksums and signed version monotonicity; the
anti-rollback barrier was exercised at each sealing (reconstruction of a
sealed version is refused, observed at the sealing gesture itself), and
version epochs open strictly monotonically --- the 0.1.1 artifact
$(0,1,1,-1,1)$ opened the new epoch; a pre-release of the sealed 0.1.0
remains refused. Sealed artifacts never changed after their sealing run
(Table~\ref{tab:empreintes}); the 0.1.1 release refreshed a stale
checksum table \emph{at birth} (17 sums, zero bytes of code moved, the
0.1.0 seal intact) --- an explicit debt-repayment pattern the protocol
names and tracks.

% ------------------------------------------------------------------
\subsection{Strengthened evidence (the v0.11.0 programme)}\label{sec:renforce}

The second edition added six laboratory instruments to the evidence base
(Table~\ref{tab:preuves110}); all run with fixed, consigned seeds and are
reproducible from the scripts shipped with this edition.

\paragraph{Mutation testing.} 832 single-byte mutants of the sealed blocks
(instructions, operands, offsets --- one byte mutated in a laboratory copy,
the sealed artifact untouched), killed by the existing instruments in
order of cost: the round-trip check, the structural oracle, and --- for
the nineteen pure-compute crypto blocks and loupe-06 --- the semantic
differential against official vectors. Result: 694/832 killed (83.4\%);
the cryptographic families reach 100\%. The 138 survivors concentrate in
the orchestration and UI blocks, where the semantic killer is not wired
at the isolated-mutant level (their semantics are verified by the full
harnesses in composition, not per mutant): this gap is \emph{named}, not
hidden, and is the primary target of the next edition.

\paragraph{Coverage-guided fuzzing.} The two untrusted-input blocks were
attacked under instrumentation: \code{loupe-01} (the file path, through
the harness's virtual OS with canaries around the path buffer) reached
60.5\,\% block-instruction coverage over 400 iterations with zero crashes
and zero canary violations; the \code{http-10} parser (on the sealed 0.4.0
artifact) reached 358 distinct instructions with zero emulation crashes
and zero named faults, growing a 17-entry corpus. The limits are stated:
coverage is per-block under emulation, and crashes are emulation faults,
not terrain deaths.

\paragraph{Bounded static memory safety.} An abstract interpretation of
all 52 blocks (the extension of the msx64 ABI audit born at run 4):
2{,}098 memory accesses classified --- 1{,}181 stack accesses verified
against the frame bound (zero out-of-frame: the stack-smash class is
absent), 263 accesses mediated by contract-bounded argument pointers, 191
constant \code{[rip]} accesses, 463 unresolved accesses \emph{named}
(volatile-recomputed bases, composed indices). The analysis is linear
(no path enumeration, loop bodies verified once, no alias analysis): a
bounded verification, not a proof.

\paragraph{Constant-time (emulation dudect).} The crypto families pass a
Welch t-test on executed-instruction counts between fixed-secret and
random-secret classes (N=60 each; all $|t| = 0$: the trace is
secret-independent --- e.g.\ x25519 executes exactly 1{,}238{,}959
instructions for every key). The metric is emulation-structural:
\emph{real-hardware wall-clock timing remains unmeasured} and is
explicitly marked as not proven.

\paragraph{Comparative benchmarks.} Against a C reference compiled from
the same contract (gcc 14.2 -O2/-Os; clang absent from the laboratory ---
the deviation is named), measured in the \emph{same} Unicorn emulator:
for SHA-256 of a 1{,}024-byte message, the hand-written blocks execute
4{,}524 instructions against the reference's 69{,}132 (-O2) and 77{,}505
(-Os) --- a 15.3$\times$ instruction advantage --- for 765 bytes of code
against 5{,}436 (-O2). The md5 and AES references were \emph{rejected by
the instrument's own self-check} (the hand-written C was faulty; the
bench refused to measure on a broken reference) and are deferred to the
next edition: better no number than a false one.

\paragraph{Oracle independence.} The oracle was re-generated by routes
independent of the code generator's: the MD5/SHA tables re-derived
through decimal arithmetic, the RFC 1321 gold vectors confronted
spec$\leftrightarrow$harness$\leftrightarrow$hashlib, the differentials
re-run with a new seed and the system generator (\code{os.urandom}):
zero divergences across all confrontations. A second \emph{agent} remains
to be tried (the replication protocol is proposed in the next edition).

% ------------------------------------------------------------------
\section{Discussion}\label{sec:discussion}

\paragraph{What the data supports.} A single developer with an LLM agent
delivered, in 17 days, a cross-platform native codebase with a pure
TLS~1.3 stack and a real Windows application, under a discipline that
made every claim checkable: every defect named, every fix diffed, every
seal fingerprinted, every decision append-only. The escaped-defect data
(RQ1/RQ3) is, to our knowledge, the first quantified account of
\emph{where} agent-written native code fails when laboratory verification
is strong but environmental fidelity is absent: not in logic, but in
loader semantics, ABI conventions and rig assumptions.

\paragraph{What the data does not support.} The study cannot claim that
the protocol is \emph{more} productive or \emph{more} correct than a
compiler-based workflow: there is no baseline, one case, one agent
family, one developer. The token budgets were measured once, early, on a
small lineage. The oracle is derived from the same layer I it checks, so
specification errors are shared, not caught. Memory safety is enforced
only to the depth of the ABI audits and the fuzzing campaigns;
constant-time behaviour of the cryptographic blocks is not analyzed in
this edition. These limits are stated, not smoothed over; the companion
instructions for the next editions (formal verification of the vault
chain, mutation testing, constant-time statistics) target them directly.

\paragraph{On the ``no compiler'' position.} We phrase the protocol's
position precisely: TRIADE eliminates the \emph{optimizing and opaque}
translation chain by having the agent write the assembly notation itself,
with a round-trip check making the notation WYSIWYG
(what-you-write-is-what-executes) at every build. It does not eliminate
the assembler, the operating system, or the platform ABI --- those are
called through their documented conventions, like any native binary.
Similarly, the sovereignty programme's claim about expropriation is
narrow: the \emph{evolution} of stolen code is blocked (signed layer-I
intent, sealed inheritance, offline keys, dead-man's-switch Shamir
escrow); the \emph{execution} of a stolen copy is not, and the
specification says so.

% ------------------------------------------------------------------
\section{Threats to validity}\label{sec:threats}

\emph{Internal.} The oracle is circular with respect to layer I
(both derive from the same specifications); a specification error passes
the oracle by construction --- the v0.11.0 independence runs (distinct
derivation routes, system generator, spec$\leftrightarrow$harness
$\leftrightarrow$hashlib tripartite confrontation) found zero
divergences, but a second \emph{agent} has not yet regenerated the
oracle. The crystallisation rule (ten identical rebuilds) checks
determinism, not correctness. The mutation-testing survivors (138/832,
concentrated in orchestration/UI blocks without an isolated-mutant
semantic killer) are a named, measured gap. \emph{External.} N=1: one
project, one developer, one agent family; no baseline comparison;
generalisation is a hypothesis for future work. \emph{Construct.} The
escaped-defect count depends on what counts as a ``named'' class --- we
counted conservatively (a class is counted once even if it killed several
runs); the token budgets measure the M1-era card, not the final
10-domain card, and measure \emph{reads}, not agent reasoning overhead.
\emph{Conclusion.} The terrain is one real Windows machine and one
developer's usage of the application; showcase liveness (0 error-reporter
events) is a weak liveness signal, not a correctness proof. Memory safety
is \emph{boundedly} verified (2{,}098 accesses classified, zero
stack-smash, 463 named unresolved --- linear analysis, no alias
analysis): a bounded result, not a proof. Constant-time is verified
\emph{structurally} (secret-independent instruction traces); real-hardware
wall-clock behaviour, cache-timing and power side channels remain
UNMEASURED and are marked as not proven.

% ------------------------------------------------------------------
\section{Related work}\label{sec:related}

\paragraph{Design by contract and specification-first development.}
TRIADE's layer I applies design-by-contract ideas~\cite{meyer1992} at the
machine-code level: one specification per block, pre/post-conditions,
named failure cases. The novelty is not the contract style but the
binding of contracts to \emph{byte ranges} through the signed card, and
the rule that no notation byte exists before its specification.

\paragraph{N-version and property-based testing.} The differential
oracles and scenario oracles follow the property-based testing
tradition~\cite{claessen2000} and the classic redundancy arguments of
N-version programming~\cite{avizienis1985}: independent instruments (three
concurrent harnesses, oracles regenerated from I, never exported) must
agree. TRIADE extends the independence requirement to the \emph{terrain}
(a machine the laboratory never controls) and to the \emph{rig itself}
(kit defects are named with the same discipline as artifact defects).

\paragraph{Certification discipline.} The milestone/exit-criteria
structure and the append-only, amend-or-invalidate decision log borrow
from airworthiness process standards such as
DO-178C~\cite{do178c}: objectives with independence, traceability, and
configuration control --- here applied to agent-written native code, with
every criterion measured rather than attested.

\paragraph{Supply-chain integrity.} Per-build fingerprints, signed
manifests, sealed tags and the key-ceremony model align TRIADE with
reproducible builds~\cite{lamb2022}, in-toto's layout
attestations~\cite{intoto2019}, SLSA's build levels~\cite{slsa},
survivable-key-compromise update systems (TUF)~\cite{tuf2010} and
automated signing at scale (Sigstore)~\cite{sigstore2022}; the offline
root-key ceremony follows the discipline of DNSSEC key
ceremonies~\cite{rfc4033}. TRIADE's addition is applying these controls
to \emph{fragments of a development state} (the card) rather than to
released artifacts only, and making the sealed tag a fact of the card
while the artifact's bytes stay frozen. Secret sharing of the master
key~\cite{shamir1979} is used both for escrow and for the dead-man's
switch.

\paragraph{Live patching.} Hot patching of native code has a long
lineage, from dynamic software updating~\cite{hicks2002} to industrial
kpatch-style trampolines~\cite{kpatch}. TRIADE's contribution is the
governance wrapper: a patch is a surgical, diffed, card-tracked revision
of blocks at stable addresses, with crystallisation, rollback and
cross-ISA rules (x86-64 trampoline vs.\ RISC-V scratch+\code{fence.i}),
validated by oracle replay and, since M9, requiring a terrain run
recouped on pieces before any seal.

\paragraph{Repository-scale context for LLM agents.} The token-budget
programme (fragmented card, one fragment per operation, no full-binary
reads) addresses the measured context limits of LLM agents: performance
degradation in long contexts~\cite{liu2024} and the need for
repository-scale retrieval and navigation~\cite{zhang2023, jimenez2024}.
The card is, in effect, a \emph{signed repo-map for machine code}: a
compact index that lets an agent operate on a large binary without ever
loading it.

\paragraph{Positioning.} None of the cited primitives is novel. The
claimed contribution is the integration --- contract-first notation at
the byte level, signed-intent governance, oracle/terrain epistemology
with named-defect accounting --- evaluated end-to-end by building a real
application with an LLM agent and reporting where it failed.

% ------------------------------------------------------------------
\section{Conclusion}\label{sec:conclusion}

We reported the first academic edition of TRIADE, a protocol under which
an LLM agent writes native code by hand under signed-intent governance,
and a longitudinal single-case evaluation of its eleven-milestone
validation programme. The data, drawn entirely from the protocol's own
append-only records, quantifies the cost (51/719/1{,}088 tokens per
locate/edit/add), the escape dynamics (eleven named defect classes over
twelve field iterations, all environmental), and the patch safety (zero
post-patch regression, armed anti-rollback) of the approach. The Loupe
application --- a real Windows hex editor born entirely from the method,
sealed twice, and alive on its user's machine --- is the existence proof
that the pipeline can carry a real deliverable. The second edition strengthened the
evidence: 83.4\% mutation kill with the cryptographic families at
100\%, zero-crash coverage-guided fuzzing of the untrusted-input blocks,
bounded memory-safety verification with zero out-of-frame accesses,
secret-independent instruction traces on the crypto families, a 15.3$\times$
instruction advantage over a compiled reference for SHA-256, and zero
divergences in oracle-independent regeneration. What remains open is now
precisely measured: the 138 mutation survivors in orchestration blocks,
the emulation-only scope of the constant-time result, the unmeasured
real-hardware timing, the missing second-agent replication --- and the
formal verification of the vault chain, which the next edition
addresses.

% ------------------------------------------------------------------
\section*{Declarations}

\paragraph{Ethics and dual-use.} The artifacts developed under TRIADE
(Loupe, the proxy, the supervisor) are the project's own tools, tested on
the developer's machines; the protocol includes no anti-analysis,
stealth or EDR-evasion mechanism, and its specification places furtivity
explicitly out of scope. The work is a dual-use-respectable engineering
study: the same discipline that hardens an artifact against tampering
also makes it auditable. The repository is published under the Business
Source License 1.1 (non-production use until the change date, then
GPL~v3).

\paragraph{Data availability.} All primary data is public in the project
repository (append-only journal, decision log, evidence vault indexes,
extraction script, generated tables). A versioned archival snapshot
(Zenodo DOI) will be minted at submission time and referenced here; the
two submission archives of this edition (\code{TRIADE\_soumission\_v0.10.0.zip}
and \code{TRIADE\_complet\_v0.10.0.zip}) contain respectively the
manuscript with its data tables, and the complete sources with raw
extracts and fingerprints.

\paragraph{Use of AI in this work.} An LLM agent (Z.ai GLM family)
executed the laboratory work of the project under the developer's
instruction, and drafted editorial material under the developer's
direction; all results reported were produced by the protocol's own
instrumentation, and the primary sources (journal, decision log,
evidence vault) were consigned session by session as the work happened.
The agent is not an author. The developer takes full responsibility for
the content of this manuscript.

\paragraph{Acknowledgements.} The author thanks the Z.ai GLM agent
(``the laboratory'') for eleven milestones of disciplined execution, and
the TRIADE decision log for keeping everyone honest.

% ------------------------------------------------------------------
\bibliography{refs}

\end{document}
